\section{Stylized Facts}

This section presents two empirical regularities that motivate the model and guide the quantitative analysis.

\subsection{Stylized Fact 1: Stablecoin Usage and Parallel Market Premia}

We document that stablecoin trading volume increases with the parallel market premium, defined as the gap between the parallel exchange rate and the official exchange rate.

Intuitively, when the official exchange rate becomes more overvalued, demand for foreign currency rises. In such environments, stablecoins provide a convenient and relatively low-cost vehicle for accessing foreign currency. As a result, their usage increases with misalignment.

This stylized fact motivates the model assumption that stablecoin demand $q_s$ is increasing in the underlying fundamental $\theta$.

\subsection{Stylized Fact 2: Stablecoin Prices Track Parallel Market Exchange Rates}

We also document that stablecoin prices in local currency terms closely track parallel market (cash) exchange rates rather than official exchange rates.

In economies with exchange rate controls, official exchange rates are often fixed or tightly managed, while parallel market rates reflect underlying supply and demand conditions. Stablecoins are typically traded on decentralized platforms without direct regulatory constraints, and their prices adjust to reflect market conditions.

As a result, stablecoin prices provide a high-frequency measure of the parallel exchange rate.

This stylized fact motivates the interpretation of stablecoin prices as a public signal about exchange rate misalignment, while cash market prices provide noisy, decentralized signals.

\subsection{Implications for the Model}

Together, these stylized facts justify two key features of the model:

\begin{enumerate}
\item Stablecoin usage is increasing in exchange rate misalignment
\item Stablecoin prices serve as an informative public signal of the parallel exchange rate
\end{enumerate}

These empirical patterns provide the foundation for the endogenous information structure in the model.
